\section{Comparison above a strict singular barrier}
\label{sec:big-minimal-comparison}

Let $\theta$ now be any smooth big real $(1,1)$-form. A strict analytic barrier replaces the smaller-background construction used above. It confines a prospective negative minimum away from its pole set; a local comparison then rules out that minimum.

Fix $\rho\leq0$ with analytic singularities and smooth remainder, and decrease the Hermitian metric if necessary so that $\theta_\rho\geq\omega$. Put $Z=\{\rho=-\infty\}$ and $\Omega=X\setminus Z$. These data remain fixed through \cref{sec:nonclosed-interior-smoothness}. A potential bounded below by $\rho-C$ is locally bounded on $\Omega$, where its extended measure agrees with its bounded product.

\begin{lemma}\label{lem:big-local-comparison}
Let $B\Subset\Omega$ be a coordinate ball, and let $u,v$ be bounded $\theta$-plurisubharmonic functions on $B$. Assume
\[
 \liminf_{x\to\partial B}(u-v)(x)\geq0,
 \qquad \mathbf1_{\{u<v\}}\mathrm{MA}_\theta(u)\leq
             \mathbf1_{\{u<v\}}\mathrm{MA}_\theta(v).
\]
Then $u\geq v$ on $B$.
\end{lemma}
\begin{proof}
Fix $\varepsilon>0$ and set $v_\varepsilon=\max(u,v-\varepsilon)$. It equals $u$ near the boundary. Locality on $\{u>v-\varepsilon\}$ and ordered contact on $\{u\leq v-\varepsilon\}$ show that
\[
 \mathrm{MA}_\theta(u)\leq\mathrm{MA}_\theta(v_\varepsilon)
\]
throughout $B$: the latter set is contained in $\{u<v\}$, where the assumed inequality applies.

Let $q$ be the ordinary plurisubharmonic envelope on $B$ of $u-v_\varepsilon$. This bounded quasi-continuous obstacle is zero near the boundary and lies between $Ah$ and zero for a sufficiently large $A$, where $h<0$ is a smooth strictly plurisubharmonic defining function for the ball. Thus $q$ tends to zero at the boundary, and \eqref{eq:local-envelope-contact} shows that $(\ddc q)^n$ is carried by $D=\{q=u-v_\varepsilon\}$.

Canonical representatives give $q+v_\varepsilon\leq u$. Ordered contact, followed by the preceding measure inequality, yields
\[
 \mathbf1_D\mathrm{MA}_\theta(q+v_\varepsilon)
 \leq\mathbf1_D\mathrm{MA}_\theta(u)
 \leq\mathbf1_D\mathrm{MA}_\theta(v_\varepsilon).
\]
The positive expansion of $(\theta_{v_\varepsilon}+\ddc q)^n$ contains both $\mathrm{MA}_\theta(v_\varepsilon)$ and $(\ddc q)^n$. Thus $\mathbf1_D(\ddc q)^n=0$, and contact support gives $(\ddc q)^n=0$. The classical bounded Bedford--Taylor comparison principle with zero boundary values implies $q=0$. Consequently $u\geq v_\varepsilon$, first quasi-everywhere and then everywhere. Let $\varepsilon\downarrow0$.
\end{proof}

\begin{theorem}\label{thm:big-barrier-domination}
Suppose $u,v\in\PSH(X,\theta)$ satisfy, for constants $C_0,C_1\geq0$,
\[
 \min(u,v)\geq\rho-C_0,\qquad v\leq u+C_1.
\]
If, for some $0\leq c<1$,
\[
 \mathbf1_{\{u<v\}}\mathrm{MA}_\theta(u)\leq
 c\,\mathbf1_{\{u<v\}}\mathrm{MA}_\theta(v),
\]
then $u\geq v$. The measures need not have finite total mass.
\end{theorem}
\begin{proof}
Replace $\rho$ by $\widehat\rho=a\rho-A$, where $a>1$ is sufficiently close to one and $A$ is sufficiently large. Since
\[
 \theta+\ddc(a\rho)=a\theta_\rho+(1-a)\theta,
\]
we can arrange $\theta_{\widehat\rho}\geq\omega/2$, $\widehat\rho-v\leq-1$, and
$u-\widehat\rho\to+\infty$ when approaching $Z$. When $Z$ is empty, the last boundary condition is unnecessary. In the proof below write $r=\widehat\rho$.

Choose $0<\delta<1$ with $(1-\delta)^n>c$, and put $v_\delta=(1-\delta)v+\delta r$. Then
\[
 u-v_\delta\geq\delta(u-r)-(1-\delta)C_1.
\]
If its infimum $m$ on $\Omega$ were negative, a sequence $x_j$ approaching that infimum would have a subsequence converging to a point $x_0\in\Omega$. Choose a coordinate ball $B\Subset\Omega$ centered at $x_0$, identify it with $\{|z|<1\}$, and take
\[
 \eta(z)=b(1-|z|^2),\qquad 0<b<1,
 \qquad \omega/2+\ddc\eta\geq0.
\]
The function $w=v_\delta+\delta\eta+m$ is $\theta$-plurisubharmonic on $B$, and $\liminf_{\partial B}(u-w)\geq0$. Also
\[
 w-v=\delta(r-v+\eta)+m<0.
\]
Thus $\{u<w\}\subset\{u<v\}$, and positivity gives on this set
\[
 \mathrm{MA}_\theta(u)\leq c\mathrm{MA}_\theta(v)\leq(1-\delta)^n\mathrm{MA}_\theta(v)\leq\mathrm{MA}_\theta(w).
\]
The local comparison lemma implies $u\geq w$ on $B$. Evaluating at $x_j$ gives $u(x_j)-v_\delta(x_j)\geq m+\delta\eta(x_j)$, contrary to $\eta(x_0)>0$. We have proved $u\geq v_\delta$. Let $\delta\downarrow0$ along values with $(1-\delta)^n>c$, and extend the resulting inequality from $\Omega$ by canonical representatives.
\end{proof}

\Needspace{20\baselineskip}
\begin{corollary}\label{cor:big-barrier-orders}
Let $u,v\in\PSH(X,\theta)$ satisfy $u,v\geq\rho-C$ for some $C\geq0$ and $u\sim v$.
\begin{enumerate}
\item If $\mathrm{MA}_\theta(u)\leq c\mathrm{MA}_\theta(v)$ for $c\geq0$, then $c\geq1$.
\item If $\mathrm{MA}_\theta(u)\leq \mathrm{e}^{\lambda(u-v)}\mathrm{MA}_\theta(v)$ for $\lambda>0$, then $u\geq v$.
\end{enumerate}
In particular the positive-exponent equation has at most one solution in each such ordinary singularity type, and the normalized equation has at most one positive scalar in that type.
\end{corollary}
\begin{proof}
If $c<1$, apply the theorem to $u,v+A$ for every constant $A>0$. Its lower-barrier constant may depend on $A$, which causes no difficulty. The conclusion $u\geq v+A$ is impossible for all $A$. For the second assertion, fix $a>0$ and restrict the measure inequality to $\{u<v-a\}$. There the factor is at most $\mathrm{e}^{-\lambda a}<1$, so the theorem gives $u\geq v-a$. Let $a\downarrow0$. Apply this assertion in both orders for two exponential solutions. Apply the first assertion in both orders for two normalized solutions.
\end{proof}

\section{Subsolutions and the Gauduchon trace}
\label{sec:big-minimal-estimates}

The strict barrier also lets us construct a subsolution within a bounded distance of any reference between $\rho$ and zero. A convex combination of positive currents gives the curvature estimate. Pairing the equation with a Gauduchon metric then bounds its normalization through a linear trace.

\begin{lemma}\label{lem:big-reference-subsolution}
Fix $p>1$ and $L>0$. There is $m>0$ such that, whenever
\[
 \rho\leq\eta\leq0,\quad\eta\in\PSH(X,\theta),
 \qquad g\geq0,\quad\|g\|_{L^p(\mathrm{d}V)}\leq L,
\]
there exists $v\in\PSH(X,\theta)$ with
\[
 \eta\leq v\leq\eta+1,\qquad
 \mathrm{MA}_\theta(v)\geq m g\,\mathrm{d}V.
\]
The constant depends only on the fixed geometry, $\rho,p,L$.
\end{lemma}
\begin{proof}
Choose $1<q<p$. The functions $h=g(2-\rho)^{2n}$ have uniformly bounded $L^q$ norms, by H\"older's inequality and the finite polynomial moments of a fixed quasi-plurisubharmonic function. Apply \cref{lem:big-lp-subsolution} with the Hermitian background $\omega$. It gives $-1\leq s\leq0$, $s\in\PSH(X,\omega)$, and
\[
 \omega_s^n\geq m h\,\mathrm{d}V
\]
for a fixed $m>0$. Put $u_1=\rho+s$. Then $\theta_{u_1}\geq\omega_s$. On $\Omega$ define
\[
 t=1+\eta-u_1,\qquad v=\eta+t^{-1}.
\]
Here $1\leq t\leq2-\rho$. Apply \cref{lem:bounded-convex-pair} on relatively compact subsets of $\Omega$, with $a=\eta$, $h=u_1$ and $\chi(s)=(1-s)^{-1}$. It gives
\[
 \theta_v=(1-t^{-2})\theta_\eta+t^{-2}\theta_{u_1}
       +2t^{-3}\,\mathrm{d}(\eta-u_1)\wedge\dc(\eta-u_1)
 \geq t^{-2}\theta_{u_1}.
\]
Since $\eta\leq v\leq\eta+1$, \cref{lem:qpsh-removable-analytic} extends $v$ across $Z$, and the bounds extend by \cref{lem:qpsh-almost-everywhere-order}. On $\Omega$, positivity of the local products gives
\[
 \mathrm{MA}_\theta(v)\geq t^{-2n}\theta_{u_1}^n
       \geq(2-\rho)^{-2n}\omega_s^n
       \geq m g\,\mathrm{d}V.
\]
Both sides give zero mass to $Z$, proving the global inequality.
\end{proof}

\begin{proposition}\label{prop:big-singular-gauduchon-trace}
Let $\gamma$ be Gauduchon, write $\gamma^n=G\,\mathrm{d}V$, and put
\[
 T_\gamma=\int_X\theta\wedge\gamma^{n-1}>0.
\]
If $u\geq\rho-C$ is $\theta$-plurisubharmonic and $\mathrm{MA}_\theta(u)\geq F\,\mathrm{d}V$, $F\geq0$, then
\begin{equation}\label{eq:big-singular-gauduchon-trace}
 \int_X F^{1/n}G^{1-1/n}\,\mathrm{d}V\leq T_\gamma.
\end{equation}
Consequently, for a fixed nonzero $0\leq f\in L^p(\mathrm{d}V)$, $p>1$, solutions above this barrier of
$\mathrm{MA}_\theta(u)=\mathrm{e}^{\lambda u}f\,\mathrm{d}V$ have uniformly bounded suprema when $\lambda>0$ is fixed. If instead $v_j\in\PSH(X,\theta)$, $j\geq1$, and constants $C_j\geq0$ satisfy
\[
 \sup_Xv_j=0,\qquad
 \mathrm{MA}_\theta(v_j)=c_j \mathrm{e}^{v_j/j}f\,\mathrm{d}V,
 \qquad v_j\geq\rho-C_j,
\]
with $c_j>0$, then the constants $c_j$ are uniformly bounded above.
\end{proposition}
\begin{proof}
The linear current $\theta_u\wedge\gamma^{n-1}$ is positive and globally defined, even at $Z$. Distributional Stokes gives
\[
 \int_X\theta_u\wedge\gamma^{n-1}
 =\int_X\theta\wedge\gamma^{n-1},
\]
since $\ddc\gamma^{n-1}=0$. Applying the identity to $\rho$ proves $T_\gamma>0$.

On $\Omega$, apply \cref{lem:interior-mixed-background} to $\theta_u$ and $\gamma$ to obtain
\[
 \theta_u\wedge\gamma^{n-1}\geq
 F^{1/n}G^{1-1/n}\,\mathrm{d}V.
\]
Exhaust $\Omega$ and use positivity of the global trace measure. The complement has zero smooth volume, giving \eqref{eq:big-singular-gauduchon-trace}.

\Needspace{5\baselineskip}
For the exponential assertion write $a=\sup_Xu$ and $w=u-a$. The estimate becomes
\[
 \mathrm{e}^{\lambda a/n}\int_X\mathrm{e}^{\lambda w/n}f^{1/n}G^{1-1/n}\,\mathrm{d}V
 \leq T_\gamma.
\]
The integral has a positive lower bound over all normalized $\theta$-plurisubharmonic $w$. Otherwise normalized compactness supplies a subsequence converging almost everywhere to a genuine potential, and dominated convergence gives a strictly positive limit integral, a contradiction. We used $w\leq0$ and the integrability of $f^{1/n}$. This bounds $a$ above.

For the last assertion the same estimate contains $c_j^{1/n}$ and $\mathrm{e}^{v_j/(jn)}$. If those integrals approached zero along a subsequence with $j\to\infty$, normalized compactness and almost-everywhere convergence would give $v_j/j\to0$ almost everywhere. Dominated convergence then gives the positive limit $\int_Xf^{1/n}G^{1-1/n}\,\mathrm{d}V$. This contradiction bounds $c_j$; finitely many remaining indices do not matter. No total top-degree mass was used in either argument.
\end{proof}
